Financial foundation models (FFMs) --- masked-pretrained, encoder-only transformers over
per-account event sequences, adapted with a lightweight head --- have been proposed to
replace the hand-engineered features and gradient-boosted decision trees (GBDTs) that
dominate production fraud detection. Such models use \emph{two} attentions: one over the
fields of an event, and one over an account's own event history. We show this design is
\emph{structurally blind} to \emph{relational} fraud --- patterns that live in a shared
entity's activity \emph{across} accounts (a compromised merchant hit by many cards; a
fraud-farm IP driving many devices) --- which no per-account model can see. We add a
\emph{third}, cross-sequence attention and study when it helps. On a controlled synthetic
benchmark with ground-truth signal placement, a per-account FFM collapses to the base rate
on relational fraud while the cross-sequence encoder recovers it. Across four real datasets
we then map \emph{when} it helps: the relational signal is sometimes a \emph{pattern} in a
neighbour's raw events (where attention over neighbours wins) and sometimes a \emph{count}
--- a velocity magnitude a scalar aggregate captures losslessly but a bounded attention
window truncates. We reconcile the two with a \emph{count-aware} cross-sequence encoder that
carries both a raw-neighbour attention path and a log-count magnitude readout. A single such
module is the best arm in \emph{both} regimes over three seeds: on real click fraud (TalkingData)
it beats the per-account FFM by \num{+0.084} PR-AUC ($0.700\pm0.011$ vs.\ $0.616\pm0.013$), and on
the synthetic benchmark it \emph{stabilises} an otherwise wildly seed-unstable raw-attention arm
($0.318\pm0.010$ vs.\ $0.252\pm0.117$). We close with the productization
levers this enables --- neighbour-embedding caching for a frozen backbone, and the count vs.\
pattern taxonomy that tells practitioners which relational encoder a given fraud problem
needs.
